\section{Introduction}
Chest X-ray (CXR) imaging is widely used to assess respiratory conditions, but interpretation is
difficult when workloads are high or specialist expertise is limited. Deep convolutional neural networks
have therefore been investigated for the automated classification of findings such as pneumonia,
tuberculosis and COVID-19 \cite{rajpurkar2017chexnet,wang2017chestxray8,jaeger2014two,cohen2020covid}.
A useful automated system must do more than score well on a held-out sample: it has to accept an image
from a user, reject unsuitable input, show what its prediction is based on, and be assessed across
disease classes and on images from sources that differ from the training data \cite{zech2018variable}.

This paper presents the design and development of LungAI, a deep-learning system for six-class CXR
classification. Given one frontal radiograph, the system predicts one of six labels: COVID-19, Normal,
Pleural Effusion, Pneumonia, Pulmonary Nodule/Mass or Tuberculosis. The classifier is an
ImageNet-initialised DenseNet-121 with a task-specific head. Around it we built a data pipeline that
assembles public collections under documented label rules, and a web application that serves the model
with an input check, Grad-CAM overlays \cite{selvaraju2017gradcam}, a low-confidence ``inconclusive''
outcome and a record of each prediction.

The contributions are:
(i) the design of an end-to-end system, from data acquisition to a served prediction, in which training,
evaluation and deployment share one preprocessing implementation;
(ii) its implementation as a tested web application with input checking, visual explanation and
safeguards against showing a prediction the model did not compute; and
(iii) an empirical assessment of the classifier on an internal test split, over three training seeds,
and on a small external cohort, including the limitations that the external evaluation reveals.
We do not introduce a new neural-network architecture and do not claim clinical readiness or
state-of-the-art performance.
